% Chinese package and font setup for Tao Slides Beamer.
% Microsoft YaHei mirrors the PowerPoint template. SimSun is a compile fallback.

\usepackage{fontspec}
\usefonttheme{professionalfonts}
\IfFileExists{/mnt/c/Windows/Fonts/msyh.ttc}{
  \setsansfont[
    Path=/mnt/c/Windows/Fonts/,
    BoldFont=msyhbd.ttc
  ]{msyh.ttc}
  \setmainfont[
    Path=/mnt/c/Windows/Fonts/,
    BoldFont=msyhbd.ttc
  ]{msyh.ttc}
}{
  \IfFontExistsTF{Microsoft YaHei}{
    \setsansfont{Microsoft YaHei}
    \setmainfont{Microsoft YaHei}
  }{
  \setsansfont{TeX Gyre Heros}
  \setmainfont{TeX Gyre Heros}
  }
}

\usepackage{xeCJK}
\IfFileExists{/mnt/c/Windows/Fonts/msyh.ttc}{
  \setCJKmainfont[
    Path=/mnt/c/Windows/Fonts/,
    BoldFont=msyhbd.ttc
  ]{msyh.ttc}
  \setCJKsansfont[
    Path=/mnt/c/Windows/Fonts/,
    BoldFont=msyhbd.ttc
  ]{msyh.ttc}
  \setCJKmonofont[
    Path=/mnt/c/Windows/Fonts/
  ]{msyh.ttc}
}{
  \IfFontExistsTF{Microsoft YaHei}{
    \setCJKmainfont[AutoFakeBold=2]{Microsoft YaHei}
    \setCJKsansfont[AutoFakeBold=2]{Microsoft YaHei}
    \setCJKmonofont{Microsoft YaHei}
  }{
  \setCJKmainfont[AutoFakeBold=2]{SimSun}
  \setCJKsansfont[AutoFakeBold=2]{SimSun}
  \setCJKmonofont{SimSun}
  }
}

\usepackage{styles/beamerthemetao}

\newcommand{\makepart}[1]{
  \part{#1}
  {
    \setbeamertemplate{footline}{}
    \begin{frame}
      \partpage
    \end{frame}
  }
}

\RequirePackage{booktabs}
\RequirePackage{colortbl}
\RequirePackage{ragged2e}
\RequirePackage{tabularx}
\RequirePackage{array}
\RequirePackage{multirow}
\RequirePackage{caption}
\RequirePackage{subcaption}
\RequirePackage{wrapfig}
\RequirePackage{pgfplots}
\RequirePackage{graphicx}
\RequirePackage{adjustbox}
\RequirePackage{environ}
\RequirePackage{textcomp}
\RequirePackage{amsmath}
\RequirePackage{amsthm}
\RequirePackage{mathtools}
\newcolumntype{Y}{>{\centering\arraybackslash}X}
\pgfplotsset{compat=1.18}

\DeclareMathOperator*{\argmax}{arg\,max}
\DeclareMathOperator*{\argmin}{arg\,min}
\setbeamertemplate{theorems}[numbered]

\theoremstyle{definition}
\newtheorem{fact}{事实}[section]
\newtheorem{examp}{示例}[section]

\theoremstyle{plain}
\newtheorem{definition}{定义}[section]
\newtheorem{proposition}{命题}
\newtheorem{theorem}{定理}
\newtheorem{assumption}{假设}

\providecommand{\H}{\mathscr{H}}
\providecommand{\E}{\mathbb{E}}
